\documentclass[publication-type=technical-report,theme=operator]{reportkit}
% Canonical fictional fixture for the experimental Operator theme.
% All product names, prices, observations, and code below are invented.
% Regenerate the analytical exhibits with:
%   PYTHONPATH=python_scripts python3 latex_templates/examples/operator-report/figures.py
\setreportkitleftheader{Workbench evaluation}
\setreportkitfooter{Fictional systems review}
\setreportkitversion{Fixture 1.0}
\setreportkitsubject{A fictional review of local-first workbench tools}
\setreportkitkeywords{ReportKit, operator theme, technical report, fictional}
\setreportkitlanguage{en-US}
\title{A workbench that stays inspectable}
\author{Systems Methods Desk}
\date{September 2026}

\begin{document}
\maketitle

\section{Decision summary}
This fictional evaluation compares three small workbench products for teams
that need local execution, visible usage costs, and a reviewable record of
changes. The exercise is deliberately narrow: it asks whether an operator can
reproduce a run, understand its limits, and explain its cost without relying
on a vendor dashboard.

\productavatar{K}{Kestrel Loom}{fictional local workbench} is the reference
candidate in this report. The compact identity \productchip{}{Morrow Dock}{hosted
runner} appears in tables, where a text label is easier to scan than an avatar.
The fictional Kestrel Team offer is \pricepill{\$18}{per seat/month}, with
\statusdot{verified} a documented local mode. A second offer,
\pricepill{\$31}{per seat/month}, adds a managed queue. These invented list
prices are scenario inputs, not market quotations.

The evidence points to a conditional recommendation. Kestrel Loom has the
clearest local audit trail in the constructed comparison, while Morrow Dock
reduces setup work for a team that already accepts hosted execution. Neither
score should be read as a quality ranking: the capability grid records only
what this fixture says is documented or not established.

\begin{decisionpoint}{Adopt a local pilot with a hosted fallback}
Run Kestrel Loom with two synthetic services for four weeks. Keep Morrow Dock
available for jobs that require a managed queue, and require an exported
run record before expanding either path.
\end{decisionpoint}

\clearpage
\section{Products and price envelope}
The comparison treats the products as separate operating models. Kestrel Loom
keeps its worktree and run ledger on the operator's machine. Morrow Dock
starts jobs in a hosted workspace. Cinder Relay accepts a local worker but
stores the queue and audit events in a fictional managed control plane. The
three names denote invented products created only for this fixture.

\begin{center}
\captionof{table}{Fictional product packages}\label{tab:operator-packages}
\begin{tabularx}{\linewidth}{@{}>{\raggedright\arraybackslash}p{3.0cm}X>{\raggedleft\arraybackslash}p{2.7cm}c@{}}
\toprule
Product & Execution and support boundary & List price & Record state \\
\midrule
\productchip{K}{Kestrel Loom}{local} & Local runner; community updates; exportable signed run ledger & \$18 / seat-month & \statusdot{verified} \\
\productchip{M}{Morrow Dock}{hosted} & Managed worker pool; email support; workspace retention for 30 days & \$24 / seat-month & \statusdot{neutral} \\
\productchip{C}{Cinder Relay}{hybrid} & Local worker with hosted queue; priority support; 90-day event history & \$31 / seat-month & \statusdot{flag} \\
\bottomrule
\end{tabularx}
\end{center}

The package rows put price beside the boundary that gives it meaning. The
managed queue is not a like-for-like substitute for a local ledger, and the
longer event history is not equivalent to a signed export. The fictional
packages are included to exercise the theme's table typography and semantic
identity markers, not to imply an external purchasing recommendation.

\begin{assumption}{Prices remain constant during the pilot}
The cost illustrations use the monthly prices shown above and do not include
tax, currency conversion, annual discounts, or support overages. A real
procurement decision would replace each value with a dated written offer.
\end{assumption}

\clearpage
\section{Capability evidence}
The grid uses a compact legend: a filled mark means the capability is
documented in this fictional evidence set; a hollow mark means it is not
established here; a dash means the capability is outside the comparison.
An unestablished mark is not a negative finding. It signals a question for
the pilot owner to close with a reproducible demonstration.

\begin{capabilitygrid}{Local,Remote,Plans,MCP,Agents,Worktree,Batch,Audit,Hooks}
\capabilityrow[fictional]{Kestrel Loom}{DDDDDDDDD}
\capabilityrow[fictional]{Morrow Dock}{UDDDDUDDU}
\capabilityrow[fictional]{Cinder Relay}{DUDDDDDDU}
\end{capabilitygrid}

For this table, “audit” means that an operator can export a record with the
input revision, selected model, tool calls, and outcome. “Hooks” means a
documented callback can stop a run before an external write. The labels avoid
claims about security certification or general product quality. The marks
\capdocumented and \capunestablished show the two evidence states in the grid.

\begin{center}
\captionof{table}{Fictional deployment envelope}\label{tab:operator-envelope}
\begin{tabular*}{\linewidth}{@{\extracolsep{\fill}}lrrrrrr@{}}
\toprule
Product & Setup (min) & Run cap & Retention (days) & Export & Queue & Offline \\
\midrule
Kestrel Loom & 14 & 80 & 365 & Signed & Local & Yes \\
Morrow Dock & 5 & 64 & 30 & JSON & Hosted & No \\
Cinder Relay & 11 & 72 & 90 & JSON & Hosted & Partial \\
\bottomrule
\end{tabular*}
\end{center}

\begin{limitationnote}{Documentation is the boundary of this comparison}
The grid reflects only the fictional briefs prepared for this example. It
does not certify that any implementation works as described, nor does it
convert missing documentation into evidence of absence.
\end{limitationnote}

\clearpage
\section{Evidence status and measured cost}
The first figure summarizes how many of the nine capability statements have
a documented result in the fixture's evidence packet. It is a count of
records, not a weighted score; every capability receives equal treatment.

\setcounter{figure}{2}
\begin{figure}[ht]
\centering
\includegraphics[width=0.91\linewidth]{figures/fig-3-evidence-status.pdf}
\caption{Documented capability statements in a fictional evidence packet.}
\label{fig:operator-evidence-status}
\end{figure}

The cost model separates one-time review work from recurring inference. In
the fictional run set, review time falls as a team reuses its runbook, while
inference cost grows with attempts. This distinction matters because a
monthly seat price alone does not reveal which operating activity drives the
total.

\begin{center}
\captionof{table}{Illustrative monthly run assumptions}\label{tab:operator-cost-inputs}
\begin{tabular}{@{}lrrr@{}}
\toprule
Scenario & Attempts & Review minutes & Inference credits \\
\midrule
Small pilot & 20 & 180 & 12 \\
Routine month & 48 & 240 & 41 \\
Credit-cap month & 80 & 270 & 64 \\
\bottomrule
\end{tabular}
\end{center}

\begin{decisionpoint}{Keep review effort visible in the operating budget}
The pilot owner should report review time and inference credits as separate
lines. A single combined cost can hide whether automation reduced operator
work or only moved it into a meter.
\end{decisionpoint}

\clearpage
\section{Cost curves and controls}
The stacked exhibit shows the fictional review and inference components for
four workloads. Hatching marks review effort; the solid segment shows
inference charges. The small final segment remains visible so a low-volume
case is not mistaken for zero consumption.

\setcounter{figure}{4}
\begin{figure}[ht]
\centering
\includegraphics[width=0.88\linewidth]{figures/fig-5-cost-per-attempt.pdf}
\caption{Fictional monthly review and inference cost by workload.}
\label{fig:operator-cost-stacks}
\end{figure}

The line exhibit holds the seat price constant and varies the number of
attempts. The vertical marker denotes a fictional credit allowance of 64
attempts; beyond it, the model charges direct inference at the same unit
rate. Solid and dashed strokes distinguish the three scenarios without
depending on hue alone.

\setcounter{figure}{5}
\begin{figure}[ht]
\centering
\includegraphics[width=0.88\linewidth]{figures/fig-6-conditional-cost.pdf}
\caption{Conditional monthly cost under three fictional run policies.}
\label{fig:operator-conditional-cost}
\end{figure}

\begin{redflag}{A credit allowance is not a spending limit}
The illustrative allowance changes the billing curve but does not stop a
run. A real pilot should use a separate attempt limit and alert before the
credit threshold is reached.
\end{redflag}

\clearpage
\section{Execution record}
An operator needs a concise record that can be checked against the revision
under review. The terminal excerpt below is invented and demonstrates the
command, comment, and output line forms. Its placeholder paths do not point
to a real system.

\begin{terminalblock}[Fictional rehearsal — workspace K-17]
\termcomment{Dry run only; no external writes are enabled.}
\termprompt{loom inspect --workspace K-17 --revision 7ac31e2}
\termline{Workspace: K-17 (fictional)}
\termline{Revision: 7ac31e2}
\termline{Policy: local-only, hooks-required}
\termline{Tools proposed: 4; tools executed: 0}
\termline{Ledger: artifacts/run-0042.json (synthetic)}
\termcomment{Review the patch before approving an execution token.}
\end{terminalblock}

The change below illustrates how a configuration review can preserve a
default-deny policy while documenting a deliberate allowance. It is long by
design: the excerpt should cross a page boundary while remaining a single
breakable diff surface. Every identifier, path, and value is fictional.

\begin{diffblock}[config/pilot-policy.yaml — fictional patch]
\diffhunk{@@ -1,18 +1,24 @@}
\diffctx{policy:}
\diffdel{  external-writes: allow}
\diffadd{  external-writes: deny}
\diffctx{  require-operator-review: true}
\diffctx{  max-attempts: 64}
\diffctx{}
\diffhunk{@@ -8,6 +14,12 @@ tools}
\diffctx{  - name: inspect-workspace}
\diffctx{    access: read-only}
\diffctx{    timeout-seconds: 20}
\diffctx{    record-result: true}
\diffctx{}
\diffadd{  - name: propose-patch}
\diffadd{    access: staged-write}
\diffadd{    approval: required}
\diffadd{    timeout-seconds: 45}
\diffadd{    record-result: true}
\diffctx{}
\diffctx{  - name: compare-run}
\diffctx{    access: read-only}
\diffctx{    timeout-seconds: 15}
\diffctx{    record-result: true}
\diffctx{}
\diffhunk{@@ -19,5 +31,11 @@ ledger}
\diffctx{ledger:}
\diffctx{  enabled: true}
\diffctx{  include-input-digest: true}
\diffctx{  include-tool-arguments: true}
\diffctx{  redact-secrets: true}
\diffctx{  retention-days: 365}
\diffctx{}
\diffadd{  export-format: jsonl}
\diffadd{  export-after-each-run: true}
\diffctx{}
\diffhunk{@@ -28,4 +46,12 @@ hooks}
\diffctx{hooks:}
\diffctx{  before-external-write:}
\diffctx{    require-named-approver: true}
\diffctx{    require-ticket-reference: true}
\diffctx{    reject-unresolved-findings: true}
\diffctx{}
\diffadd{  before-model-retry:}
\diffadd{    max-retries: 2}
\diffadd{    preserve-previous-output: true}
\diffadd{    record-reason: true}
\diffctx{}
\diffhunk{@@ -35,3 +61,14 @@ reporting}
\diffctx{reporting:}
\diffctx{  summary: concise}
\diffctx{  include-cost-estimate: true}
\diffctx{  include-policy-decisions: true}
\diffctx{  include-revision: true}
\diffctx{}
\diffadd{  evidence-bundle:}
\diffadd{    include-run-ledger: true}
\diffadd{    include-patch-digest: true}
\diffadd{    include-tool-inventory: true}
\diffadd{    include-approvals: true}
\diffadd{    include-limitations: true}
\diffctx{}
\diffhunk{@@ -41,2 +78,8 @@ workspace}
\diffctx{workspace:}
\diffctx{  retain-on-failure: true}
\diffadd{  cleanup-requires-review: true}
\diffadd{  preserve-failed-run: true}
\diffadd{  export-path: artifacts/pilot-run.jsonl}
\diffctx{}
\diffhunk{@@ -45,1 +86,5 @@ operator}
\diffctx{operator:}
\diffctx{  alias: sample-reviewer}
\diffadd{  second-review-required: true}
\diffadd{  approval-expires-hours: 8}
\diffctx{  acknowledge-limits: true}
\diffctx{}
\diffhunk{@@ -48,0 +94,7 @@}
\diffadd{audit:}
\diffadd{  schema-version: 1}
\diffadd{  checksum-algorithm: sha256}
\diffadd{  verify-before-export: true}
\diffadd{  include-clock-source: true}
\diffadd{  include-runtime-version: true}
\diffadd{  include-policy-version: true}
\diffctx{}
\diffhunk{@@ -52,0 +103,8 @@}
\diffadd{limits:}
\diffadd{  max-output-bytes: 64000}
\diffadd{  max-tool-calls: 24}
\diffadd{  max-wall-seconds: 900}
\diffadd{  stop-on-missing-evidence: true}
\diffadd{  stop-on-policy-mismatch: true}
\diffadd{  stop-on-unknown-tool: true}
\diffctx{}
\diffhunk{@@ -96,0 +170,6 @@}
\diffadd{closing:}
\diffadd{  result: no-external-changes}
\diffadd{  ledger: synthetic-example}
\diffadd{  approval: not-requested}
\diffadd{  next-step: inspect-rendered-pages}
\diffctx{}
\end{diffblock}

\clearpage
\section{Pilot boundary}
The patch is illustrative, not a deployment template. Before a real trial,
the operator would replace its synthetic identity, validate every hook
against the selected runtime, set an owner for data retention, and rehearse
both the denied-write path and the recovery path. The report's purpose is to
make those obligations visible beside the product comparison and cost model.

\begin{evidencenote}{What this fixture demonstrates}
The document uses semantic product identity, a nine-capability grid, named
statuses, price pills, four existing callout types, three booktabs tables,
three chart APIs, and line-based terminal and diff surfaces. Every measured
value belongs to the fictional scenario declared in the neighboring figure
script.
\end{evidencenote}

The pilot should end with a short operator record: the exact revision,
approved scope, attempted actions, rejected actions, export checksum, and
open questions. If that record cannot be reconstructed by another reviewer,
the system is not ready for a wider trial.

\end{document}
